%!TEX root = ../../thesis.tex

\section{Reinforcement Learning}\label{sec:reinforcement-learning}

% Drafting outline: bullet points to be fleshed out into prose.

\subsection{Agent, Environment, and Markov Decision Processes}\label{sec:rl-mdp}

% --- Verbatim excerpt from the proposal (own text; citation keys remapped) ---
Reinforcement learning (RL) is an extremely general framework for approaching problems
in which an agent strives to maximize a reward signal through a selection of actions
in certain states~\cite{sutton_reinforcement_2020}.
These kinds of problems can be modeled as finite Markov decision processes (MDPs).
Formally, a Markov decision process is a 3-tuple $MDP=(\mathcal{S},\mathcal{A},p)$ where

\begin{itemize}
  \item $\mathcal{S}$ is the set of states,
  \item $\mathcal{A}$ is the set of actions where $\mathcal{A}_s\subseteq A$ are actions available in state $s\in\mathcal{S}$,
  \item $p(s^\prime,r|s,a)=\Pr\{S_t=s^\prime, R_t=r|S_{t-1}=s, A_{t-1}=a\}$
    describes the system dynamics, i.e., the probability of encountering a next
    state $s^\prime$ and reward $r$ at time step $t$ when taking action $a$ in
    state $s$ at time step $t-1$. This satisfies the \emph{Markov property}:
    $p$ depends only on the current state and not the full history.
\end{itemize}

An RL agent's goal is to learn a policy $\pi:\mathcal{S}\times\mathcal{A}\rightarrow [0,1]$
mapping each state $s$ to probabilities of taking each possible action $a$ to maximize
the cumulative reward.\\

A common way to visualize MDPs is via a directed graph: nodes model the state
space $S$ while edges $(s_i,s_j)$ labeled with an action-reward-probability
tuple $(a,r,p)\in\mathcal{A}\times\mathcal{R}\times[0,1]$ model potential
transitions.

A standard example introduced in~\textcite{sutton_reinforcement_2020} is that of a \emph{recycling robot} (\cref{fig:mdp-recycling-robot}).
The robot collects empty cans and, at each decision point, either actively searches
for one, waits for someone to hand it a can, or returns to its charging station.
Its state is the charge level of its battery, $\mathcal{S}=\{\mathtt{high},\mathtt{low}\}$.

%!TEX root = ../thesis.tex
% Second example MDP for \cref{sec:rl-mdp}: the recycling robot of Sutton and Barto
% (Example 3.3), redrawn in the (a,r,p) edge convention of \cref{fig:mdp-example}.
\begin{figure}[htbp]
  \centering
  \begin{tikzpicture}[
      >={Stealth[round]},
      every state/.style={draw,minimum size=1.1cm,inner sep=1pt},
      every edge/.style={draw,->,font=\small},
      every loop/.style={->,looseness=8},
    ]
    \node[state] (high) at (0,0)   {$\mathtt{high}$};
    \node[state] (low)  at (7.6,0) {$\mathtt{low}$};

    \path
      (high) edge[loop above] node[above] {$(\mathtt{search},r_{\mathrm{s}},\alpha)$}      (high)
      (high) edge[loop below] node[below] {$(\mathtt{wait},r_{\mathrm{w}},1)$}             (high)
      (low)  edge[loop above] node[above] {$(\mathtt{search},r_{\mathrm{s}},\beta)$}       (low)
      (low)  edge[loop below] node[below] {$(\mathtt{wait},r_{\mathrm{w}},1)$}             (low)
      (high) edge[bend left=32] node[above] {$(\mathtt{search},r_{\mathrm{s}},1-\alpha)$}  (low)
      (low)  edge node[above] {$(\mathtt{recharge},0,1)$}                                  (high)
      (low)  edge[bend left=32] node[below] {$(\mathtt{search},-3,1-\beta)$}               (high);
  \end{tikzpicture}
  \caption[Example MDP: the recycling robot]{%
    The recycling robot of Sutton and Barto~\cite{sutton_reinforcement_2020}.
    Its state space consists of the charge level of its battery, $\mathcal{S}=\{\mathtt{high},\mathtt{low}\}$,
    and its action space consists of $\mathcal{A}=\{\mathtt{search},\mathtt{wait},\mathtt{recharge}\}$.
    Its battery state transition is probabilistic and depends on $\alpha$ and $\beta$.
    $r_s$ and $r_w$ are the rewards for search and waiting.
  }
  \label{fig:mdp-recycling-robot}
\end{figure}


A simple greedy approach is to always select the $\mathtt{search}$ action since it
pays the highest rewards in expectation.
This approach encurs the penalty of searching while the battery is in the $\mathtt{low}$.
An agent must therefore recognize the state he is in and choose its action accordingly.
While this example is easily solved, MDPs with larger state and action spaces present
a much harder problem.

\subsubsection{Partially Observable Markov Decision Processes}\label{sec:rl-pomdp}

In the previous formulation of MDPs, the true state of the environment is available
to the agent.
In reality, this is often not the case.
An agent perceives its environment as a lossy function of the underlying state.
The recycling robot of \cref{fig:mdp-recycling-robot}, for instance, might perceive
its charge level as a noisy voltage measurement.
Observing the environment may therefore introduce uncertainty.
The Markov property then still holds for the underlying state, but not for what the agent sees.

This situation is modeled through a \emph{partially observable Markov decision process}
(POMDP)~\cite{ASTROM1965174,kaelbling_planning_1998}, which
extends the original 3-tuple MDP $(\mathcal{S},\mathcal{A},p)$ to a 5-tuple POMDP=$(\mathcal{S},\mathcal{A},p,\Omega,o)$, by adding

\begin{itemize}
  \item $\Omega$, the set of observations the agent can receive, and
  \item $o(\omega|s^\prime,a)=\Pr\{O_t=\omega|S_t=s^\prime,A_{t-1}=a\}$, the observation
    function, giving the probability of perceiving $\omega$ after a transition into
    $s^\prime$ under action $a$.
\end{itemize}

The environment still evolves as an MDP over $\mathcal{S}$ according to $p$.
However, the agent must act on observation $O_t$ instead of on the true state $S_t$.
An ordinary MDP is the special case $\Omega=\mathcal{S}$ with $o=\text{id}(\cdot)$.

POMDPs introduce the problem of \emph{state aliasing}.
If two states that require different actions map to the same observation, then no deterministic
memoryless policy $\pi:\Omega\times\mathcal{A}\rightarrow[0,1]$ can act correctly in both,
and the best achievable memoryless policy can be arbitrarily worse than a policy that is
allowed to consult the history~\cite{SINGH1994284}.

A POMDP can be made Markovian by treating not the actual observations, but the
entire observation-action history $H_t=O_0,A_0,O_1,A_1,\dots,O_t$ as the state space.
Another classical solution is to summarize the observation-action history in a probability distribution 
and treat this distribution as the actual state.
This state is called a \emph{belief state} $b_t(s)=\Pr\{S_t=s|H_t\}$ and suffices
to allow the agent to act optimally~\cite{puterman_markov_2009}.
Practically speaking, the belief-state approach is restricted in that it presupposes
the existence of a known model $(p,o)$. Solving the resulting belief MDP is also intractable
for all but the simplest problems.

This reduction is mostly of theoretical value here, since it presupposes a known model
$(p,o)$ and an explicit state space, and solving the resulting belief MDP exactly is
intractable for all but very small problems~\cite{papadimitriou_complexity_1987}.

\subsection{The Exploration-Exploitation Tradeoff}
